\clearpage
\section{Analiza problemu}
\label{sec:analiza}

\subsection{Zastosowanie i reprezentacja danych}

Drzewo BST służy do porządkowania danych według klucza.
Porównanie z wartością bieżącego węzła pozwala wybrać jedno
z poddrzew, zamiast sprawdzać wszystkie elementy.
Struktura ta może być wykorzystywana do przechowywania zbiorów
unikatowych wartości i budowy prostych indeksów.

W projekcie węzeł \kod{Node} zawiera wartość \kod{value}
oraz pola \kod{left} i \kod{right}. Pola dzieci są indeksami
w wektorze \kod{tree}; wartość \kod{-1} oznacza brak dziecka.
Indeks korzenia znajduje się w \kod{rootIndex}.
Dla pustego drzewa również wynosi on \kod{-1}.

Rozmieszczenie węzłów w wektorze jest oddzielone od ich połączeń.
Element może leżeć pod dowolnym poprawnym indeksem, o ile pola
dzieci prowadzą do właściwych poddrzew. Sama kolejność wektora
nie musi być posortowana.

\subsection{Sposób działania algorytmu}

Dodawanie rozpoczyna się od korzenia. Wartość mniejsza jest kierowana
do lewego poddrzewa, a większa do prawego. Po dojściu do pustego
połączenia nowy węzeł zostaje dopisany na końcu wektora,
a jego indeks zapisany jako połączenie od rodzica.
Jeżeli wartość już występuje, operacja nie dodaje duplikatu.

Usuwanie rozpoczyna się od wyszukania wartości i zapamiętania
rodzica znalezionego węzła. Gdy wartości nie ma, drzewo pozostaje
bez zmian. Dla znalezionego elementu rozpatrywane są trzy przypadki:
\begin{itemize}
  \item liść zostaje odłączony od rodzica;
  \item węzeł z jednym dzieckiem zostaje zastąpiony połączeniem
        do tego dziecka;
  \item węzeł z dwojgiem dzieci otrzymuje wartość minimum prawego
        poddrzewa, po czym usuwany jest węzeł tego następnika.
\end{itemize}

Usunięcie pozycji wektora przesuwa późniejsze elementy.
Program zmniejsza więc o jeden wszystkie indeksy większe od
usuniętej pozycji, uwzględniając również indeks korzenia.
Metoda \kod{removeTree()} czyści cały wektor i ustawia
\kod{rootIndex} na \kod{-1}. Obiekt może następnie przyjmować
kolejne wartości.

Dla wysokości drzewa $h$ przejście po jednej ścieżce ma koszt $O(h)$.
Dodawanie, wyszukiwanie wartości i wyszukiwanie ścieżki zaczynają się
od korzenia i korzystają z tej reguły. Drzewo nie jest równoważone,
więc przy niekorzystnej kolejności danych
wysokość może rosnąć liniowo z liczbą węzłów $n$.
Skuteczne usunięcie ma koszt $O(n)$ ze względu na przesuwanie
elementów i poprawianie indeksów. Realokacja bufora wektora
może również nadać pojedynczemu dodaniu koszt $O(n)$.

\subsection{Ręczny przykład działania}
\label{sec:przyklad-reczny}

Po dodaniu wartości $1,2,3,8,5$ korzeniem jest 1.
Węzły 2, 3 i 8 tworzą ciąg prawych połączeń, natomiast
5 jest lewym dzieckiem węzła 8.
Rysunek~\ref{rys:przyklad-bst} pokazuje tę strukturę
oraz wynik usunięcia wartości 3. Litery L i P oznaczają
lewe i prawe dziecko.

\begin{figure}[H]
\centering
\begin{tikzpicture}[
  wezel/.style={circle,draw,minimum size=6.5mm,inner sep=1pt,font=\small},
  krawedz/.style={-{Stealth[length=1.6mm]},thin},
  etykieta/.style={font=\scriptsize,fill=white,inner sep=1pt}
]
\node[font=\small\bfseries] at (0.9,0.8) {Przed usunięciem 3};
\node[font=\small\bfseries] at (6.5,0.8) {Po usunięciu 3};

\node[wezel] (a1) at (0,0) {1};
\node[wezel] (a2) at (0.9,-0.9) {2};
\node[wezel] (a3) at (1.8,-1.8) {3};
\node[wezel] (a8) at (2.7,-2.7) {8};
\node[wezel] (a5) at (1.8,-3.6) {5};
\draw[krawedz] (a1) -- node[etykieta,above right] {P} (a2);
\draw[krawedz] (a2) -- node[etykieta,above right] {P} (a3);
\draw[krawedz] (a3) -- node[etykieta,above right] {P} (a8);
\draw[krawedz] (a8) -- node[etykieta,above left] {L} (a5);

\node[wezel] (b1) at (5.6,0) {1};
\node[wezel] (b2) at (6.5,-0.9) {2};
\node[wezel] (b8) at (7.4,-1.8) {8};
\node[wezel] (b5) at (6.5,-2.7) {5};
\draw[krawedz] (b1) -- node[etykieta,above right] {P} (b2);
\draw[krawedz] (b2) -- node[etykieta,above right] {P} (b8);
\draw[krawedz] (b8) -- node[etykieta,above left] {L} (b5);
\end{tikzpicture}

\caption{Drzewo przed usunięciem wartości 3 i po jej usunięciu}
\label{rys:przyklad-bst}
\end{figure}

Węzeł o wartości 3 ma jedno dziecko: węzeł 8.
Program łączy więc rodzica, czyli węzeł 2, bezpośrednio z węzłem 8.
Następnie usuwa pozycję o indeksie 2 z wektora.
Węzły 8 i 5 przesuwają się pod indeksy 2 i 3,
dlatego trzeba poprawić prowadzące do nich połączenia.

Tabela~\ref{tab:zmiana-indeksow} przedstawia zawartość wektora
przed operacją i po korekcie indeksów. Korzeń nadal ma indeks 0.
Kreski w ostatnim wierszu oznaczają pozycję, której po usunięciu
nie ma już w kontenerze.

\begin{table}[H]
\centering
\caption{Zmiana zawartości wektora po usunięciu wartości 3}
\label{tab:zmiana-indeksow}
\small
\begin{tabular}{|c|c|c|c|c|c|c|}
\hline
& \multicolumn{3}{c|}{\textbf{Przed usunięciem}}
& \multicolumn{3}{c|}{\textbf{Po usunięciu}}\\ \hline
\textbf{Indeks} & \textbf{value} & \textbf{left} & \textbf{right}
& \textbf{value} & \textbf{left} & \textbf{right}\\ \hline
0 & 1 & $-1$ & 1 & 1 & $-1$ & 1\\ \hline
1 & 2 & $-1$ & 2 & 2 & $-1$ & 2\\ \hline
2 & 3 & $-1$ & 3 & 8 & 3 & $-1$\\ \hline
3 & 8 & 4 & $-1$ & 5 & $-1$ & $-1$\\ \hline
4 & 5 & $-1$ & $-1$ & --- & --- & ---\\ \hline
\end{tabular}
\end{table}

Po tej operacji metoda \kod{show()} wypisałaby $1,2,8,5$,
zgodnie z kolejnością w wektorze. Opcja inorder w menu podąża
za połączeniami i wypisuje wartości w kolejności $1,2,5,8$.

\subsection{Zależności między modułami a kolejność scalania}
\label{sec:kolejnosc-scalania}

Jeżeli obie gałęzie robocze wprowadzają zależne zmiany, najpierw
należy scalić do \kod{main} operacje drzewa. Na przykład zapis
binarny używa metody \kod{preorderValues()}, a menu korzysta
z \kod{findPath()}. Zmiany w deklaracjach tych metod powinny
trafić do gałęzi z interfejsem, zanim jej PR zostanie scalony.
W tym celu włącza się do niej aktualny stan \kod{main}
zgodnie z procedurą z punktu~\ref{sec:konflikt-scalania}.

Po aktualizacji gałęzi interfejsu sprawdza się kompilację programu
i uruchamia próby zależnych operacji. Dla \kod{findPath()} można
porównać ścieżkę do znanego węzła, a dla zapisu binarnego
zapisać drzewo, opróżnić je i odtworzyć z pliku. Wyniki
tych prób można dołączyć do PR-a jako podstawę przeglądu.
